\lecture{4}{Sun 02 Feb 2025 16:14}{I had the flu and missed this class :(}
\begin{proposition}\label{prp:1.4}
	If $X$ is Hausdorff, then every singleton $\left\{ x \right\} $ is closed.
\end{proposition}
\begin{proof}
	It suffices to show $\left\{ x \right\}^{c} $ is open. Since $X$ is Hausdorff, for all $y\neq x$, there exist disjoint neighborhoods $U_x,U_y$ of $x,y$, respectively. Let $\left\{ U_\alpha  \right\}_{\alpha \in X\setminus \left\{ x \right\} } $ be a collection of neighborhoods such that for all $\alpha $, there exists an open set $V_\alpha  $ containing $x$ such that $U_\alpha \cap V_\alpha =\varnothing $. It follows that

	\[
		X\setminus \left\{ x \right\} =\bigcup_{\alpha \in X\setminus \left\{ x \right\} } U_\alpha 
	.\]
	Therefore, $\left\{ x \right\} ^{c} $ is closed.
\end{proof}
\chapter{Interior, Closure, and Boundary}
We want to be able to talk more precisely about the interior, exterior, and boundary of a topological space.
\begin{definition}[Interior]\label{dfn:2.1}
	Let $X$ be a topological space, and let $A\subseteq X$. The \emph{interior} of $A$, denoted $Int(A)$, is the union of all open sets contained in $A$. It is the largest open set contained in $A$.
\end{definition}
\begin{definition}[Closure]\label{dfn:2.2}
	Let $X$ be a topological space, and let $A\subseteq X$. The \emph{closure} of $A$, denoted $\overline{A} $, is the intersection of all closed sets containing $A$. It is the smallest closed set containing $A$.
\end{definition}
Note that by definition, $\Int(A)\subseteq A\subseteq \overline{A} $. This fact will be useful on multiple occasions.
\begin{proposition}\label{prp:2.1}
	Let $X$ be a topological space, and let $A\subseteq X$. Then
	\begin{itemize}
		\item $A$ is open if and only if $\Int(A)=A$;
		\item $A$ is closed if and only if $\overline{A} =A$.
	\end{itemize}
\end{proposition}
\begin{proof}
	Suppose $A$ is open. Then since $A$ contains itself, and $\Int(A)$ is a union over all open sets contained in $A$, we have $A\subseteq \Int(A)$. By definition, $\Int(A)\subseteq A$. Now suppose $A=\Int(A)$. Since $\Int(A)$ is a union of open sets, it must be open, and thus $A $ is open.

	Suppose $A$ is closed. Then since $A$ contains itself, and $\overline{A} $ is an intersection over all closed sets containing $A$, we have $\overline{A} \subseteq A$. By definition, $A\subseteq \overline{A} $. Now suppose $A=\overline{A} $. By proposition \ref{prp:1.3}, since $\overline{A} $ is an intersection of closed sets, it is closed, and thus $A $ is closed.
\end{proof}

Here's an interesting example. Consider the finite compliment topology. Closed sets are, by definition, finite, with the sole exception of $X$. Assuming $X$ is infinite, then $\overline{A} =X$ for all \emph{open} sets $A$.

Here's another, very important example. Consider $\mathbb{Q} \subseteq \mathbb{R} $ with the standard topology. Note that $\Int(\mathbb{Q} )=\varnothing $, because for any open interval $(a,b)$, there will be an irrational number contained within. Similarly, $\overline{\mathbb{Q} } =\mathbb{R}  $.
\begin{definition}[Dense]\label{dfn:2.3}
	A subset $A\subseteq X$ of a topological space $X$ is called \emph{dense} if $\overline{A} =X$.
\end{definition}
\begin{proposition}\label{prp:2.2}
	Let $A\subseteq X$, and let $X$ be a topological space. Then
	\begin{itemize}
		\item $\Int(A^{c} )=\left( \overline{A}  \right) ^{c} $;
		\item $\overline{\left( A^{c}  \right) } =\Int(A)^{c} $.
	\end{itemize}
\end{proposition}
\begin{proof}
	The argument is basically the same for both implications, so I will only do the first for convenience. Let $\mathcal{A} $ be the collection of all open sets contained within $A$, and recall that
	\[
		\Int(A)^{c} =\left( \bigcup_{U\in \mathcal{A} } U \right) ^{c} 
	.\]
	By De Morgan's law, this equals
	\[
		\bigcap_{U\in \mathcal{A} } U^{c} 
	.\]
	Since a set is closed if and only if it's the compliment of an open set, this collects all closed sets containing $A^{c} $. Therefore, this is $\left( \overline{A}  \right) ^{c} $.
\end{proof}
\begin{proposition}\label{prp:2.3}
	Let $X$ be a topological space, and let $A\subseteq X$. Then $x\in \Int(A)$ if and only if there exists an open neighborhood $U\subseteq A$ of $x$. Additionally, $y\in \overline{A} $ if and only if every neighborhood of $y$ intersects $A$.
\end{proposition}
\begin{proof}
	The first proposition is trivial. $\Int(A)$ is open by definition, so an open neighborhood exists. If an open neighborhood $U\subseteq A$ exists, then $x\in \Int(A)$ by definition.

	For the second proposition, suppose $y\in \overline{A} $, and suppose for contradiction there exists some open set $U$ containing $y$ such that $U \cap A=\varnothing $. It follows that $U\subseteq A^{c} $, and thus $y\in \Int(A^{c} )$ by the first proposition. By proposition \ref{prp:2.2}, $\Int(A^{c} )^{c} =\overline{A} $, and thus we have $y\notin \overline{A} $, a contradiction.

	Now suppose for all neighborhoods $U$ of $y$ that $U \cap A$ is nonempty. Suppose for contradiction that $y\notin \overline{A} $. It follows that $y\in \overline{A} ^{c} $, and by proposition \ref{prp:2.2}, $y\in \Int(A^{c} ) $. By the first proposition, there exists an open set $U\subseteq A^{c} $, and thus $U \cap A=\varnothing $, a contradiction.
\end{proof}
This proposition motivates the following definition.
